\chapter{Transformada de Fourier de tiempo continuo}

\section{Introducción}

\begin{enumerate}
	
	\item Vamos a ver como podemos \textbf{modificar} el concepto de \textbf{serie de Fourier} \textbf{para estudiar funciones no periódicas} de tiempo continuo. Esto nos permite \textbf{definir} el concepto de \textbf{contenido de frecuencias de una señal}, aún cuando sea no periódica.
	
	\item Pensemos en el ejemplo de un tren de pulsos simétrico, donde
	\begin{equation}
		a_{0} = \frac{T_{1}}{T_{2}},
	\end{equation}
	y
	\begin{equation}
		a_{k} = \frac{T_{1}}{T} \, \frac{\sin\left(\frac{k \, \pi \, T_{1}}{T}\right)}{\frac{k \, \pi \, T_{1}}{T}}. \label{tf.1}
	\end{equation}
	
	\item Podemos escribir (\ref{tf.1}) de la siguiente forma
	\begin{equation}
		T a_{k} = \left. T_{1} \, \frac{\sin(\omega)}{\omega} \right|_{\omega = \frac{k \, \pi \, T_{1}}{T}}
	\end{equation}
	y entonces es evidente que la podemos interpretar como una \textbf{secuencia} formada por \textbf{muestras} de una \textbf{función envolvente} $T_{1} \, \frac{\sin(\omega)}{\omega}$ donde la variable $\omega$ es continua, y la secuencia $T a_{k}$ está formada por valores igualmente espaciados en frecuencia, con intervalo $\frac{2 \, \pi}{T}$ (suponiendo $2T_1$ como el ancho del pulso).
	
	\item Observe que la separación de las muestras de la envolvente de $T a_{k}$ está dada por 
	\begin{equation}
		\omega_{T} = \frac{2 \, \pi}{T}.
	\end{equation}
	Entonces, cuanto mayor sea $T$ menor será $\omega_{T}$ y la cantidad de muestras en un intervalo dado de frecuencia será mayor.
	
	\item Si fijamos el valor del ancho del pulso (recordando que debe ser menor a $T$, por hipótesis) la envolvente de $T a_{k}$ es independiente de $T$.
	
\end{enumerate}

\section{Síntesis y análisis}

\subsection{Teoría}

\begin{enumerate}
	
	\item Vamos a construir una señal $x(t)$ no periódica como el límite de una señal $\hat{x}(t)$ periódica cuando el período $T$ tiende a $\infty$. Veamos qué pasa con los coeficientes de la serie de Fourier. Recordemos que una señal periódica $\hat{x}(t) = \hat{x}(t+T)$ se puede escribir como una serie de Fourier mediante la ecuación de síntesis 
	\begin{equation}
		\hat{x}(t) = \sum_{k=-\infty}^{\infty} a_{k} \, e^{j \, k \, \omega_{0} \, t},
	\end{equation}
	donde los coeficientes $a_{k}$ se calculan mediante la ecuación de análisis 
	\begin{equation}
		a_{k} = \frac{1}{T} \int_{-\frac{T}{2}}^{\frac{T}{2}} \hat{x}(t) e^{-j \, k \, \omega_{0} \, t} \, dt.
	\end{equation}
	
	\item Definamos la señal no periódica $x(t)$ especificando 
	\begin{equation}
		x(t) = \begin{cases}
			\hat{x}(t), & |t| \leq \frac{T}{2}, \\
			0, & |t| > \frac{T}{2},
		\end{cases}
		\qquad
		\hat{x}(t) = \hat{x}(t+T) .
	\end{equation}
	Entonces los coeficientes de la serie de Fourier son 
	\begin{equation}
		a_{k} = \frac{1}{T} \int_{-\frac{T}{2}}^{\frac{T}{2}} \hat{x}(t) e^{-j \, k \, \omega_{0} \, t} \, dt = \frac{1}{T} \int_{-\infty}^{\infty} x(t) e^{-j \, k \, \omega_{0} \, t} \, dt.
	\end{equation}
	\textbf{Las dos integrales son iguales en valor porque un período de $\hat{x}(t)$ coincide con $x(t),$ justamente donde se evalúa la primer integral, y $x(t)$ es nula para $|t| > \frac{T}{2}.$}
	
	\item Definamos la ecuación de análisis de la transformada de Fourier como 
	\begin{equation}
		X(j \, \omega) = \int_{-\infty}^{\infty} x(t) e^{-j \, \omega \, t} \, dt.
	\end{equation}
	Tenemos que los coeficientes de la serie de Fourier forman la secuencia de valores 
	\begin{equation}
		T a_{k} = X(j \, k \, \omega_{0}) .
	\end{equation}
	
	\item Usando la ecuación de síntesis de la serie de Fourier podemos escribir 
	\begin{equation}
		\hat{x}(t) = \frac{1}{T} \sum_{k=-\infty}^{\infty} X(j \, k \, \omega_{0}) e^{j \, k \, \omega_{0} \, t}.
	\end{equation}
	Dado que $\omega_{0} = \frac{2 \, \pi}{T} \longleftrightarrow T = \frac{2 \, \pi}{\omega_{0}}$ vemos que la serie de Fourier puede escribirse como
	\begin{equation}
		\hat{x}(t) = \frac{1}{2 \, \pi} \sum_{k=-\infty}^{\infty} X(j \, k \, \omega_{0}) e^{j \, k \, \omega_{0} \, t} \, \omega_{0}.
	\end{equation}
	A medida que $T \rightarrow \infty$, $\omega_{0} \rightarrow 0$ (o sea un diferencial de frecuencia $d\omega$), $\hat{x}(t) \rightarrow x(t)$ y la sumatoria de Riemann se convierte en integral, obteniendo la ecuación de síntesis:
	\begin{equation}
		x(t) = \frac{1}{2 \, \pi} \int_{-\infty}^{\infty} X(j \, \omega) e^{j \, \omega \, t} \, d\omega .
	\end{equation}
	\textbf{A medida que $T \rightarrow \infty$, se agranda el rango de valores de $t$ en el que las dos funciones $\hat{x}(t)$ y $x(t)$ son iguales.}
\end{enumerate}

\subsection{Prestar atención}

Como hemos visto antes, vamos a usar la ecuación de análisis 
\begin{equation}
	X(j \, \omega) = \int_{-\infty}^{\infty} x(t) e^{-j \, \omega \, t} \, dt,  \label{analisis}
\end{equation}
y la ecuación de síntesis 
\begin{equation}
	x(t) = \frac{1}{2 \, \pi} \int_{-\infty}^{\infty} X(j \, \omega) e^{j \, \omega \, t} \, d\omega . \label{sintesis}
\end{equation}
Esto significa que estamos usando la variable $\omega$ (radianes por segundo), y no la variable $f$ (ciclos por segundo o Hertz). Recordemos que ambas están relacionadas por 
\begin{equation}
	\omega = 2 \, \pi f.
\end{equation}
Este cambio de variables afecta a las ecuaciones que serían de la forma 
\begin{equation}
	X(f) = \int_{-\infty}^{\infty} x(t) e^{-j 2 \, \pi f t} \, dt,
\end{equation}
y 
\begin{equation}
	x(t) = \int_{-\infty}^{\infty} X(f) e^{j 2 \, \pi f t} \, df.
\end{equation}
Usamos las fórmulas dadas por las ecs. (\ref{analisis}) y (\ref{sintesis}) porque son las que aparecen en la convención de los libros estándar de la asignatura (como Oppenheim). La diferencia de notación no afecta la teoría.

\subsection{Concepto de la Transformada de Fourier}

\begin{enumerate}
	
	\item La Transformada de Fourier es una generalización de la serie de Fourier, y permite estudiar funciones no periódicas en el dominio de la frecuencia.
	\item La Transformada de Fourier permite descomponer una señal en componentes armónicas cuyas frecuencias varían en diferenciales, a diferencia de la serie de Fourier donde las frecuencias de las armónicas varían en incrementos discretos de $\omega_{0} = \frac{2 \, \pi}{T}$, la frecuencia fundamental.
\end{enumerate}

\section{Convergencia de la Transformada de Fourier}

Sea una señal $x(t),$ y consideremos las funciones $X(j \, \omega)$ y $\hat{x}(t)$ (o $x_F(t)$) definidas como 
\begin{equation}
	X(j \, \omega) = \int_{-\infty}^{\infty} x(t) e^{-j \, \omega \, t} \, dt,  \label{convergenciaUno}
\end{equation}
y 
\begin{equation}
	x_{F}(t) = \frac{1}{2 \, \pi} \int_{-\infty}^{\infty} X(j \, \omega) e^{j \, \omega \, t} \, d\omega .
\end{equation}
La pregunta es: ¿cuándo $x_{F}(t)$ es una representación válida de $x(t)$?

\subsection{Energía finita}

Si $x(t)$ tiene energía finita 
\begin{equation}
	\int_{-\infty}^{\infty} |x(t)|^{2} \, dt < \infty ,
\end{equation}
entonces $X(j \, \omega)$ es finita, o sea que la integral (\ref{convergenciaUno}) converge, y además 
\begin{equation}
	\int_{-\infty}^{\infty} |x(t) - x_{F}(t)|^{2} \, dt = 0.
\end{equation}
Esto implica que $x_{F}(t)$ y $x(t)$ son similares excepto en, a lo sumo, una cantidad limitada de valores de $t$ donde las funciones tengan discontinuidades.

\subsection{Condiciones de Dirichlet}

Son \textbf{condiciones alternativas} a la de energía, y son \textbf{suficientes} para asegurar que $x_{F}(t)$ es igual a $x(t)$ excepto en puntos de discontinuidad.
\begin{enumerate}
	\item $x(t)$ debe ser absolutamente integrable:
	\begin{equation}
		\int_{-\infty}^{\infty} |x(t)| \, dt < \infty .
	\end{equation}
	
	\item $x(t)$ tiene un número finito de máximos y mínimos dentro de cualquier intervalo finito.
	\item $x(t)$ tiene un número finito de discontinuidades dentro de cualquier intervalo finito, y dichas discontinuidades son finitas. En este caso, en los puntos de discontinuidad de $x(t)$, la transformada inversa $x_{F}(t)$ es igual al promedio de los valores de $x(t)$ a ambos lados de la discontinuidad.
\end{enumerate}
Si se cumplen estas condiciones, entonces $x(t)$ tiene transformada de Fourier.

\subsection{Pregunta}

Las señales que llegan a las antenas de un televisor o un equipo de radio, por mencionar dos ejemplos, ¿tienen transformada de Fourier? Justifique.

\section{Ejemplos}

\subsection{Impulso temporal}

Calculemos la transformada de Fourier de un impulso $\delta(t)$
\begin{equation}
	X(j \, \omega) = \int_{-\infty}^{\infty} \delta(t) e^{-j \, \omega \, t} \, dt.
\end{equation}
Por la propiedad del muestreo 
\begin{equation}
	\delta(t) e^{-j \, \omega \, t} = \delta(t) e^{0} = \delta(t) ,
\end{equation}
y por lo tanto 
\begin{equation}
	X(j \, \omega) = \int_{-\infty}^{\infty} \delta(t) e^{-j \, \omega \, t} \, dt = \int_{-\infty}^{\infty} \delta(t) \, dt = 1.
\end{equation}
Note que es un valor constante para todo valor de $\omega.$

\subsection{Impulso temporal desplazado}

Calculemos la transformada de Fourier de un impulso desplazado en el tiempo $\delta(t - t_{0})$
\begin{equation}
	X(j \, \omega) = \int_{-\infty}^{\infty} \delta(t - t_{0}) e^{-j \, \omega \, t} \, dt.
\end{equation}
Recordemos que, por la propiedad del muestreo,
\begin{equation}
	\delta(t - t_{0}) e^{-j \, \omega \, t} = \delta(t - t_{0}) e^{-j \, \omega \, t_{0}}
\end{equation}
y por lo tanto podemos escribir
\begin{equation}
	X(j \, \omega) = \int_{-\infty}^{\infty} \delta(t - t_{0}) e^{-j \, \omega \, t_{0}} \, dt = e^{-j \, \omega \, t_{0}} \int_{-\infty}^{\infty} \delta(t - t_0) \, dt = e^{-j \, \omega \, t_{0}},
\end{equation}
teniendo en cuenta que el factor $e^{-j \, \omega \, t_{0}}$ no depende de la variable de integración $t$ y puede ser extraído de la integral como una constante.

\subsection{Impulso en frecuencia}

Sea la transformada de Fourier 
\begin{equation}
	X(j \, \omega) = 2 \, \pi \, \delta(\omega - \omega_{0}) .
\end{equation}
Apliquemos la ecuación de síntesis para obtener 
\begin{equation}
	x(t) = \frac{1}{2 \, \pi} \int_{-\infty}^{\infty} 2 \, \pi \, \delta(\omega - \omega_{0}) \, e^{j \, \omega \, t} \, d\omega = e^{j \, \omega_{0} \, t}.
\end{equation}
Esto significa que las funciones periódicas exponenciales complejas $e^{j \, \omega_{0} \, t}$ también tienen transformada de Fourier, y son impulsos en frecuencia $2 \, \pi \, \delta(\omega - \omega_{0})$.

\subsection{Tren de impulsos en frecuencia}

En general, si tenemos un tren de impulsos en frecuencia
\begin{equation}
	X(j \, \omega) = \sum_{k=-\infty}^{\infty} 2 \, \pi \, a_{k} \, \delta(\omega - k \, \omega_{0}) ,
\end{equation}
al aplicar la inversa obtenemos la serie de Fourier:
\begin{equation}
	x(t) = \sum_{k=-\infty}^{\infty} a_{k} \, e^{j \, k \, \omega_{0} \, t},
\end{equation}
por lo que podemos escribir la dualidad:
\begin{equation}
	\sum_{k=-\infty}^{\infty} a_{k} \, e^{j \, k \, \omega_{0} \, t} 
	\stackrel{\mathcal{F}}{\longleftrightarrow}
	\sum_{k=-\infty}^{\infty} 2 \, \pi \, a_{k} \, \delta(\omega - k \omega_{0}). \label{tf.tif}
\end{equation}

\subsection{Tren de impulsos temporal}

Consideremos el tren de impulsos temporal
\begin{equation}
	x(t) = \sum_{k=-\infty}^{\infty} \delta(t - kT) .
\end{equation}
Calculemos los coeficientes de la serie de Fourier, y observemos que dentro del intervalo fundamental de integración (de $-T/2$ a $T/2$) tenemos $x(t) = \delta(t)$, lo que permite escribir:
\begin{equation}
	a_{k} = \frac{1}{T} \int_{-\frac{T}{2}}^{\frac{T}{2}} \delta(t) e^{-j \, k \, \omega_{0} \, t} \, dt = \frac{1}{T}.
\end{equation}
Y por lo tanto, aplicando la propiedad (\ref{tf.tif}), tenemos que la transformada de Fourier de un tren de impulsos en el tiempo es otro tren de impulsos en frecuencia:
\begin{equation}
	X(j \, \omega) = \sum_{k=-\infty}^{\infty} \frac{2 \, \pi}{T} \, \delta\left(\omega - \frac{2 \, \pi \, k}{T}\right).
\end{equation}

\subsection{Pulso rectangular en tiempo}

Calculemos la transformada de la señal pulso rectangular 
\begin{equation}
	x(t) = \begin{cases}
		1, & |t| \leq T_{1}, \\ 
		0, & |t| > T_{1}.
	\end{cases}
\end{equation}
Según la ecuación de análisis 
\begin{equation}
	X(j \, \omega) = \int_{-T_{1}}^{T_{1}} e^{-j \, \omega \, t} \, dt = \frac{\left[e^{-j \, \omega \, t}\right]_{-T_{1}}^{T_{1}}}{-j \, \omega}.
\end{equation}
Evaluando la integral
\begin{equation}
	X(j \, \omega) = \frac{e^{-j \, \omega \, T_{1}} - e^{j \, \omega \, T_{1}}}{-j \, \omega}.
\end{equation}
Finalmente, multiplicando y diviendo por $2$ (o $2T_1$) y usando la identidad de Euler para el seno llegamos a la expresión:
\begin{equation}
	X(j \, \omega) = 2 \, \frac{e^{j \, \omega \, T_{1}} - e^{-j \, \omega \, T_{1}}}{2j \, \omega} = 2 \, \frac{\sin(\omega \, T_{1})}{\omega}.
\end{equation}
Que puede escribirse elegantemente como:
\begin{equation}
	X(j \, \omega) = 2 \, T_{1} \, \frac{\sin(\omega \, T_{1})}{\omega \, T_{1}}.
\end{equation}

\subsection{Pulso rectangular en frecuencia}

Veamos cuál es la señal $x(t)$ cuya transformada de Fourier es un filtro pasabajos ideal:
\begin{equation}
	X(j \, \omega) = \begin{cases}
		1, & |\omega| \leq W_{1}, \\ 
		0, & |\omega| > W_{1}.
	\end{cases}
\end{equation}
Según la ecuación de síntesis 
\begin{equation}
	x(t) = \frac{1}{2 \, \pi} \int_{-W_{1}}^{W_{1}} e^{j \, \omega \, t} \, d\omega = \frac{1}{2 \, \pi} \, \frac{\left[e^{j \, \omega \, t}\right]_{-W_{1}}^{W_{1}}}{j \, t}.
\end{equation}
Evaluando la integral
\begin{equation}
	x(t) = \frac{1}{2 \, \pi} \, \frac{e^{j \, W_{1} \, t} - e^{-j \, W_{1} \, t}}{j \, t}.
\end{equation}
Reordenando los términos para que coincidan con la definición del seno:
\begin{equation}
	x(t) = \frac{1}{\pi t} \left( \frac{e^{j \, W_{1} \, t} - e^{-j \, W_{1} \, t}}{2j} \right) = \frac{\sin(W_{1} \, t)}{\pi \, t}.
\end{equation}
Por lo tanto, resulta la clásica función Sinc:
\begin{equation}
	x(t) = \frac{W_{1}}{\pi} \frac{\sin(W_{1} \, t)}{W_{1} \, t}.
\end{equation}

\subsection{Exponencial decreciente derecha}

Consideremos la señal causal
\begin{equation}
	x(t) = e^{-a \, t} \, u(t), \quad a > 0.
\end{equation}
Según la ecuación de análisis 
\begin{equation}
	X(j \, \omega) = \int_{0}^{\infty} e^{-a \, t} \, e^{-j \, \omega \, t} \, dt = \left[ \frac{e^{-(a + j \, \omega) t}}{-(a + j \, \omega)} \right]_{0}^{\infty}.
\end{equation}
Dado que $a>0$, el límite en infinito converge a cero, por lo que:
\begin{equation}
	X(j \, \omega) = \frac{1}{a + j \, \omega}, \quad a > 0.
\end{equation}
Observe además que la señal es absolutamente integrable:
\begin{equation}
	\int_{0}^{\infty} | e^{-a \, t} \, u(t) | \, dt = \frac{1}{a}.
\end{equation}
y la función $x(t)$ tiene una única discontinuidad acotada en $t = 0$. Por lo tanto, esta función verifica plenamente las condiciones de Dirichlet.

\section{Propiedades de la Transformada de Fourier de tiempo continuo}

\subsection{Linealidad}

Sea una función combinada:
\begin{equation}
	r(t) = a \, x(t) + b \, y(t).
\end{equation}
Podemos probar que su transformada es:
\begin{equation}
	R(j \, \omega) = \int_{-\infty}^{\infty} (a \, x(t) + b \, y(t)) e^{-j \, \omega \, t} \, dt.
\end{equation}
Por la linealidad del operador integral, esto se puede separar como:
\begin{equation}
	R(j \, \omega) = a \left( \int_{-\infty}^{\infty} x(t) e^{-j \, \omega \, t} \, dt \right) + b \left( \int_{-\infty}^{\infty} y(t) e^{-j \, \omega \, t} \, dt \right).
\end{equation}
Por lo tanto resulta:
\begin{equation}
	R(j \, \omega) = a \, X(j \, \omega) + b \, Y(j \, \omega).
\end{equation}
Así deducimos la propiedad:
\begin{equation}
	a \, x(t) + b \, y(t) \stackrel{\mathcal{F}}{\longleftrightarrow} a \, X(j \, \omega) + b \, Y(j \, \omega) .
\end{equation}

\subsection{Desplazamiento en el tiempo}

Dada la función retardada
\begin{equation}
	r(t) = x(t - t_{0}) ,
\end{equation}
su transformada es
\begin{equation}
	R(j \, \omega) = \int_{-\infty}^{\infty} x(t - t_{0}) e^{-j \, \omega \, t} \, dt.
\end{equation}

Aplicando un cambio de variable $\tau = t - t_{0}$ (donde $dt = d\tau$), obtenemos:
\begin{equation}
	R(j \, \omega) = \int_{-\infty}^{\infty} x(\tau) \, e^{-j \, \omega (\tau + t_{0})} \, d\tau.
\end{equation}
Separando el término exponencial:
\begin{equation}
	R(j \, \omega) = e^{-j \, \omega \, t_{0}} \int_{-\infty}^{\infty} x(\tau) e^{-j \, \omega \tau} \, d\tau.
\end{equation}
Así, llegamos a la conclusión:
\begin{equation}
	x(t - t_{0}) \stackrel{\mathcal{F}}{\longleftrightarrow} e^{-j \, \omega \, t_{0}} X(j \, \omega) .
\end{equation}

\subsection{Conjugación y simetría del conjugado}

\begin{enumerate}
	\item El conjugado de la transformada de Fourier es 
	\begin{equation}
		\overline{X(j \, \omega)} = \overline{\int_{-\infty}^{\infty} x(t) \, e^{-j \, \omega \, t} \, dt} = \int_{-\infty}^{\infty} \overline{x(t)} \, e^{j \, \omega \, t} \, dt.
	\end{equation}
	Cambiando el signo de $\omega$ obtenemos:
	\begin{equation}
		\overline{X(-j \, \omega)} = \int_{-\infty}^{\infty} \overline{x(t)} \, e^{-j \, \omega \, t} \, dt.
	\end{equation}
	Como conclusión:
	\begin{equation}
		\overline{x(t)} \stackrel{\mathcal{F}}{\longleftrightarrow} \overline{X(-j \, \omega)} .
	\end{equation}
	
	\item Si $x(t)$ es puramente real, entonces $x(t) = \overline{x(t)}$. Sustituyendo esto, tenemos:
	\begin{equation}
		\overline{X(-j \, \omega)} = \int_{-\infty}^{\infty} x(t) e^{-j \, \omega \, t} \, dt = X(j \, \omega),
	\end{equation}
	lo que demuestra la propiedad de simetría conjugada para señales reales:
	\begin{equation}
		X(-j \, \omega) = \overline{X(j \, \omega)} . \label{conjuuno}
	\end{equation}
	
	\item De la Ec. (\ref{conjuuno}), si $x(t)$ es real se desprende que la parte real del espectro es una función par, y la parte imaginaria es impar:
	\begin{align}
		\Re(X(j \, \omega)) &= \Re(X(-j \, \omega)), \\
		\Im(X(j \, \omega)) &= -\Im(X(-j \, \omega)).
	\end{align}
	
	\item En términos polares, si $x(t)$ es real, el módulo al cuadrado es:
	\begin{equation}
		|X(j \, \omega)|^{2} = X(j \, \omega) \, \overline{X(j \, \omega)} = X(j \, \omega) \, X(-j \, \omega) = |X(-j \, \omega)|^{2}.
	\end{equation}
	Por lo tanto el módulo $|X(j \, \omega)|$ es una función par, y la fase $\angle X(j \, \omega)$ es una función impar.
	
	\item Si $x(t)$ es real y par ($x(t) = x(-t)$), entonces $X(j \, \omega)$ es puramente real y par.
	\item Si $x(t)$ es real e impar ($x(t) = -x(-t)$), entonces $X(j \, \omega)$ es puramente imaginaria e impar.
\end{enumerate}

\subsection{Diferenciación en el tiempo}

\begin{enumerate}
	\item Escribamos la ecuación de síntesis:
	\begin{equation}
		x(t) = \frac{1}{2 \, \pi} \int_{-\infty}^{\infty} X(j \, \omega) e^{j \, \omega \, t} \, d\omega .
	\end{equation}
	
	\item Derivemos ambos términos respecto al tiempo:
	\begin{equation}
		\frac{dx(t)}{dt} = \frac{1}{2 \, \pi} \int_{-\infty}^{\infty} X(j \, \omega) \frac{d(e^{j \, \omega \, t})}{dt} \, d\omega.
	\end{equation}
	
	\item Obtenemos finalmente:
	\begin{equation}
		\frac{dx(t)}{dt} = \frac{1}{2 \, \pi} \int_{-\infty}^{\infty} j \, \omega \, X(j \, \omega) e^{j \, \omega \, t} \, d\omega.
	\end{equation}
	
	\item La conclusión es que derivar en el tiempo equivale a multiplicar por $j\omega$ en frecuencia:
	\begin{equation}
		\frac{dx(t)}{dt} \stackrel{\mathcal{F}}{\longleftrightarrow} j \, \omega \, X(j \, \omega) .
	\end{equation}
	
\end{enumerate}

\subsection{Diferenciación en frecuencia}

\begin{enumerate}
	\item Escribamos la ecuación de análisis:
	\begin{equation}
		X(j \, \omega) = \int_{-\infty}^{\infty} x(t) e^{-j \, \omega \, t} \, dt.
	\end{equation}
	
	\item Derivemos ambos términos respecto a la frecuencia:
	\begin{equation}
		\frac{dX(j \, \omega)}{d\omega} = \int_{-\infty}^{\infty} x(t) \, \frac{d(e^{-j \, \omega \, t})}{d\omega} \, dt.
	\end{equation}
	
	\item Obtenemos:
	\begin{equation}
		\frac{dX(j \, \omega)}{d\omega} = \int_{-\infty}^{\infty} -j \, t \, x(t) \, e^{-j \, \omega \, t} \, dt.
	\end{equation}
	
	\item Multiplicando por $j$ en ambos lados:
	\begin{equation}
		j \, \frac{dX(j \, \omega)}{d\omega} = \int_{-\infty}^{\infty} t \, x(t) \, e^{-j \, \omega \, t} \, dt.
	\end{equation}
	
\end{enumerate}

Conclusión:
\begin{equation}
	t \, x(t) \stackrel{\mathcal{F}}{\longleftrightarrow} j \, \frac{dX(j \, \omega)}{d\omega} .
\end{equation}


\subsection{Escalado en tiempo y frecuencia}

\begin{enumerate}
	
	\item Vamos a calcular la transformada de $x(a \, t)$ en función de la transformada de $x(t)$.
	\item Escribamos la ecuación de análisis:
	\begin{equation}
		\mathcal{F}\{x(at)\} = \int_{-\infty}^{\infty} x(a \, t) e^{-j \, \omega \, t} \, dt,
	\end{equation}
	donde $a \neq 0$. 
	
	\item Hagamos la sustitución $\tau = a \, t$ (con $dt = d\tau / a$). Debemos cuidar los límites de integración según el signo de $a$:
	\begin{equation}
		\mathcal{F}\{x(at)\} = \begin{cases}
			\frac{1}{a} \int_{-\infty}^{\infty} x(\tau) e^{-j \frac{\omega}{a} \tau} \, d\tau , & a > 0, \\ 
			-\frac{1}{a} \int_{-\infty}^{\infty} x(\tau) e^{-j \frac{\omega}{a} \tau} \, d\tau , & a < 0.
		\end{cases}
	\end{equation}
	
	\item El signo de $a$ afecta al orden de integración. Si se conservan los límites originales, debe absorberse el signo mediante el valor absoluto. La conclusión es:
	\begin{equation}
		x(at) \stackrel{\mathcal{F}}{\longleftrightarrow} \frac{1}{|a|} X\left( \frac{j \, \omega}{a} \right) ,
	\end{equation}
	donde $a \neq 0$.
	
\end{enumerate}

\subsection{Convolución en el dominio del tiempo}

\begin{enumerate}
	\item Empezamos con la integral de convolución:
	\begin{equation}
		y(t) = \int_{-\infty}^{\infty} x(\tau) \, h(t-\tau) \, d\tau .
	\end{equation}
	\item Deseamos obtener la transformada de Fourier:
	\begin{equation}
		Y(j \, \omega) = \int_{-\infty}^{\infty} \left[ \int_{-\infty}^{\infty} x(\tau) \, h(t-\tau) \, d\tau \right] e^{-j \, \omega \, t} \, dt.
	\end{equation}
	\item Cambiando el orden de integración, y observando que $x(\tau)$ no depende de $t$, obtenemos:
	\begin{equation}
		Y(j \, \omega) = \int_{-\infty}^{\infty} x(\tau) \left( \int_{-\infty}^{\infty} h(t-\tau) \, e^{-j \, \omega \, t} \, dt \right) d\tau .
	\end{equation}
	\item Por la propiedad de desplazamiento en el tiempo, la integral interior es la transformada de una señal retardada $\tau$:
	\begin{equation}
		\int_{-\infty}^{\infty} h(t-\tau) \, e^{-j \, \omega \, t} \, dt = e^{-j \, \omega \tau} \, H(j \, \omega).
	\end{equation}
	Reemplazando esto en la ecuación principal obtenemos:
	\begin{equation}
		Y(j \, \omega) = \int_{-\infty}^{\infty} x(\tau) \, e^{-j \, \omega\tau} \, H(j \, \omega) \, d\tau .
	\end{equation}
	\item Extrayendo $H(j \, \omega)$ (que no depende de $\tau$) de la integral, nos queda:
	\begin{equation}
		Y(j \, \omega) = \left( \int_{-\infty}^{\infty} x(\tau) e^{-j \, \omega \tau} \, d\tau \right) H(j \, \omega) = X(j \, \omega) \, H(j \, \omega) .
	\end{equation}
\end{enumerate}

\textbf{Conclusión fundamental: la convolución de señales en el dominio del tiempo implica la multiplicación de sus espectros en el dominio de la frecuencia.}
\begin{equation}
	x(t) * h(t) \stackrel{\mathcal{F}}{\longleftrightarrow} X(j \, \omega) \cdot H(j \, \omega) .
\end{equation}

\subsection{Convolución en el dominio de la frecuencia (Modulación)}

\begin{enumerate}
	\item Empezamos definiendo un espectro $R(j\omega)$ como la convolución de dos espectros:
	\begin{equation}
		R(j \, \omega) = \frac{1}{2 \, \pi} \int_{-\infty}^{\infty} X(j \, \theta) Y(j(\omega - \theta)) \, d\theta .
	\end{equation}
	\item Aplicando la ecuación de síntesis y manipulando el orden de integración (proceso dual al anterior), llegamos a que la señal temporal correspondiente es:
	\begin{equation}
		r(t) = y(t) \cdot x(t) .
	\end{equation}
\end{enumerate}

\textbf{Conclusión: la convolución de transformadas en frecuencia implica la multiplicación de señales en el dominio del tiempo.}
\begin{equation}
	x(t) \cdot y(t) \stackrel{\mathcal{F}}{\longleftrightarrow} \frac{1}{2 \, \pi} X(j \, \omega) * Y(j \, \omega).
\end{equation}

\subsection{Traslación en frecuencia}

Partiendo de la integral de síntesis, si multiplicamos una señal por una exponencial compleja $e^{j \, \omega_{0} t}$ (lo que altera su fase de forma lineal), equivale a un desplazamiento de su espectro en la frecuencia por $\omega_0$.
\begin{equation}
	e^{j \, \omega_{0} t} x(t) \stackrel{\mathcal{F}}{\longleftrightarrow} X(j(\omega - \omega_{0})) .
\end{equation}
Observemos la dualidad con la propiedad de desplazamiento temporal:
\begin{equation}
	x(t - t_{0}) \stackrel{\mathcal{F}}{\longleftrightarrow} e^{-j \, \omega \, t_{0}} X(j \, \omega) .
\end{equation}

\subsection{Relación de Parseval}

La energía total de una señal puede calcularse tanto en el dominio del tiempo como integrando el cuadrado de su espectro en el dominio de la frecuencia:
\begin{equation}
	\int_{-\infty}^{\infty} |x(t)|^{2} \, dt = \frac{1}{2 \, \pi} \int_{-\infty}^{\infty} |X(j \, \omega)|^{2} \, d\omega .
\end{equation}

\subsection{Dualidad}

La similitud estructural entre las ecuaciones de análisis y síntesis da lugar a la propiedad de dualidad. Si una señal $x(t)$ tiene un espectro $X(j\omega)$ (o $X(f)$), entonces si generamos una señal temporal con la "forma" de $X(t)$, su espectro tendrá la forma de $x(-\omega)$. En términos de la frecuencia $f$ en Hertz:
\begin{align}
	x(t) &\stackrel{\mathcal{F}}{\longleftrightarrow} X(f), \\
	X(t) &\stackrel{\mathcal{F}}{\longleftrightarrow} x(-f).
\end{align}

\subsection{Integración}

Dada una señal $x(t)$, su integral es la función $\int_{-\infty}^{t} x(v) \, dv$. Se puede demostrar que integrar en el tiempo equivale a dividir por $j\omega$ en frecuencia, añadiendo un término de impulso en $\omega=0$ que representa el valor medio o el área acumulada de la señal:
\begin{equation}
	\int_{-\infty}^{t} x(\tau) \, d\tau \stackrel{\mathcal{F}}{\longleftrightarrow} \frac{1}{j \, \omega} X(j \, \omega) + \pi X(0) \delta(\omega).
\end{equation}

Para probar esta propiedad, debemos encontrar primero la transformada del escalón $u(t)$.

\subsubsection{Escalón $u(t)$}

\begin{enumerate}
	
	\item Primero encontremos la transformada de una señal simétrica 
	\begin{equation}
		x(t) = -\frac{1}{2} + u(t).
	\end{equation}
	
	\item Dado que la señal $x(t)$ es constante en $-\frac{1}{2}$ y $\frac{1}{2}$ excepto por un salto de amplitud $1$ en $t=0$, la derivada de $x(t)$ es igual al impulso unitario $\delta(t)$.
	Entonces, por la propiedad de la derivada, $j\omega X(j\omega) = 1$, por lo que la transformada de $x(t)$ debe ser:
	\begin{equation}
		-\frac{1}{2} + u(t) \stackrel{\mathcal{F}}{\longleftrightarrow} \frac{1}{j \omega}. \label{fourierintegraluno}
	\end{equation}
	
	\item Recordemos que la transformada de Fourier de una constante $K$ es $2\pi K \delta(\omega)$, por ende la de la señal constante $\frac{1}{2}$ es $\pi \delta(\omega)$. Despejando $u(t)$ de la Ec. (\ref{fourierintegraluno}) y empleando la propiedad de linealidad, obtenemos finalmente:
	\begin{equation}
		u(t) \stackrel{\mathcal{F}}{\longleftrightarrow} \pi \delta(\omega) + \frac{1}{j \omega}.
	\end{equation}
	
\end{enumerate}

\subsubsection{Integral}

\begin{enumerate}
	
	\item La integral de $x(t)$ puede ser expresada como la convolución con un escalón:
	\begin{equation}
		\int_{-\infty}^{t} x(v) \, dv = x(t) * u(t).
	\end{equation}
	Entonces, por la propiedad de la convolución, la transformada de Fourier de la integral es la multiplicación de los espectros: $X(j\omega) \cdot U(j\omega)$.
	\item Por lo tanto:
	\begin{equation}
		\int_{-\infty}^{t} x(\tau) \, d\tau \stackrel{\mathcal{F}}{\longleftrightarrow} X(j \, \omega) \left( \frac{1}{j \, \omega} + \pi \delta(\omega) \right) = \frac{X(j\omega)}{j \omega} + \pi X(0) \delta(\omega).
	\end{equation}
\end{enumerate}
